% Options for packages loaded elsewhere
\PassOptionsToPackage{unicode}{hyperref}
\PassOptionsToPackage{hyphens}{url}
\documentclass[
  11pt,
  openany]{book}
\usepackage{xcolor}
\usepackage{amsmath,amssymb}
\setcounter{secnumdepth}{5}
\usepackage{iftex}
\ifPDFTeX
  \usepackage[T1]{fontenc}
  \usepackage[utf8]{inputenc}
  \usepackage{textcomp} % provide euro and other symbols
\else % if luatex or xetex
  \usepackage{unicode-math} % this also loads fontspec
  \defaultfontfeatures{Scale=MatchLowercase}
  \defaultfontfeatures[\rmfamily]{Ligatures=TeX,Scale=1}
\fi
\usepackage{lmodern}
\ifPDFTeX\else
  % xetex/luatex font selection
  \setmainfont[]{Times New Roman}
\fi
% Use upquote if available, for straight quotes in verbatim environments
\IfFileExists{upquote.sty}{\usepackage{upquote}}{}
\IfFileExists{microtype.sty}{% use microtype if available
  \usepackage[]{microtype}
  \UseMicrotypeSet[protrusion]{basicmath} % disable protrusion for tt fonts
}{}
\makeatletter
\@ifundefined{KOMAClassName}{% if non-KOMA class
  \IfFileExists{parskip.sty}{%
    \usepackage{parskip}
  }{% else
    \setlength{\parindent}{0pt}
    \setlength{\parskip}{6pt plus 2pt minus 1pt}}
}{% if KOMA class
  \KOMAoptions{parskip=half}}
\makeatother
\usepackage{longtable,booktabs,array}
\usepackage{calc} % for calculating minipage widths
% Correct order of tables after \paragraph or \subparagraph
\usepackage{etoolbox}
\makeatletter
\patchcmd\longtable{\par}{\if@noskipsec\mbox{}\fi\par}{}{}
\makeatother
% Allow footnotes in longtable head/foot
\IfFileExists{footnotehyper.sty}{\usepackage{footnotehyper}}{\usepackage{footnote}}
\makesavenoteenv{longtable}
\setlength{\emergencystretch}{3em} % prevent overfull lines
\providecommand{\tightlist}{%
  \setlength{\itemsep}{0pt}\setlength{\parskip}{0pt}}
\usepackage[]{natbib}
\bibliographystyle{apalike}
\usepackage{booktabs}
\usepackage[top=1cm,bottom=1cm,left=1.5cm,right=1.5cm,includehead]{geometry}
\usepackage{fancyhdr}
  \pagestyle{fancy}
  \fancypagestyle{plain}{%
  \renewcommand{\headrulewidth}{0pt}%
  \fancyhf{}%
}
  \fancyhf{}
  \fancyhf[HRO,HLE]{\thepage}
  \fancyhf[HC]{\leftmark}
\usepackage{quotchap}
\usepackage{amsmath}
\usepackage[most]{tcolorbox}
\usepackage{tikz}
\usetikzlibrary{patterns}
\usetikzlibrary{decorations.pathreplacing,angles,quotes}
\usepackage{titlesec}
\titlespacing*{\chapter}{0pt}{0cm}{0cm}
\titlespacing*{\section}{0pt}{0.2cm}{0.2cm}
\titlespacing*{\subsection}{0pt}{0.2cm}{0.2cm}
\titlespacing*{\subsubsection}{0pt}{0.2cm}{0.2cm}
\usepackage{caption}
\usepackage{subcaption}
\usepackage{units}
\usepackage{amssymb}
\usepackage{pifont}
\usepackage{nth}
\usepackage{mathtools}
\usepackage{dirtytalk}
\usepackage{setspace}
\usepackage{multicol}
\usepackage[draft]{tikzpeople}
\usepackage{hyperref}
\usepackage{gensymb}
\usepackage{tabularx}
\usepackage{multirow}
\usepackage{array}
\usepackage{wrapfig}
\graphicspath{ {images/} }
\usepackage{exsheets}
\usepackage{setspace}
\usepackage{etoolbox}
\usepackage{enumitem}
\usepackage{bm}
\usepackage{vwcol}
\usetikzlibrary{shapes.geometric}
\usepackage{graphicx}


\newcommand{\PreserveBackslash}[1]{\let\temp=\\#1\let\\=\temp}
\newcolumntype{C}[1]{>{\PreserveBackslash\centering}p{#1}}
\newcolumntype{R}[1]{>{\PreserveBackslash\raggedleft}p{#1}}
\newcolumntype{L}[1]{>{\PreserveBackslash\raggedright}p{#1}}

\newenvironment{cols}[1][]{}{}

\newenvironment{col}[1]{\begin{minipage}{#1}\ignorespaces}{%
\end{minipage}
\ifhmode\unskip\fi
\aftergroup\useignorespacesandallpars}

\def\useignorespacesandallpars#1\ignorespaces\fi{%
#1\fi\ignorespacesandallpars}

\makeatletter
\def\ignorespacesandallpars{%
  \@ifnextchar\par
    {\expandafter\ignorespacesandallpars\@gobble}%
    {}%
}
\makeatother


\tcbuselibrary{skins}

\usepackage{graphbox,graphicx}

\pretolerance=100
\tolerance=200
\hyphenchar\font=-1
\sloppy

\usepackage{enumitem}
\setlist{topsep=0pt, leftmargin=*}
\setlist[2]{itemsep=1pt}



\usepackage{float}
\usepackage{adjustbox}

\usepackage{pgfplots}
\usepackage{pgfmath}
\pgfplotsset{compat=1.8}
\pgfmathdeclarefunction{gauss}{3}{%
\pgfmathparse{1/(#3*sqrt(2*pi))*exp(-((#1-#2)^2)/(2*#3^2))}%
}

\usepackage{epigraph}
\usepackage[switch, modulo]{lineno}

\usepackage[dvipsnames]{xcolor}
\usepackage[no-math]{fontspec}
\usepackage{lmodern}
\renewcommand{\familydefault}{lmss}


\AtBeginDocument{\frontmatter}
\usepackage{bookmark}
\IfFileExists{xurl.sty}{\usepackage{xurl}}{} % add URL line breaks if available
\urlstyle{same}
\hypersetup{
  pdftitle={Higher Apps Bookdown},
  pdfauthor={A James},
  hidelinks,
  pdfcreator={LaTeX via pandoc}}

\title{Higher Apps Bookdown}
\author{A James}
\date{2026-04-15}

\begin{document}
\maketitle

{
\setcounter{tocdepth}{1}
\tableofcontents
}
\chapter*{Course Information}\label{course-information}
\addcontentsline{toc}{chapter}{Course Information}

\part{Modelling}\label{part-modelling}

\chapter{Developing Models}\label{developing-models}

\section{Fermi Estimation}\label{fermi-estimation}

\section{Graphing Relationships}\label{graphing-relationships}

\section{Linear Models}\label{linear-models}

\section{Quadratic Models}\label{quadratic-models}

\section{Exponential Models}\label{exponential-models}

\section{Recurrence Models}\label{recurrence-models}

\section*{Review}\label{review}
\addcontentsline{toc}{section}{Review}

\chapter{Units of Measure}\label{units-of-measure}

\section{Test}\label{test}

\section{Test}\label{test-1}

\section{Test}\label{test-2}

\section{Test}\label{test-3}

\subsection{Testing}\label{testing}

\subsection{Testing}\label{testing-1}

\section{Test}\label{test-4}

\subsection{Testing}\label{testing-2}

\subsubsection{Still Testing}\label{still-testing}

\chapter{Error and Tolerance}\label{error-and-tolerance}

\chapter{Using Models}\label{using-models}

\chapter{Spreadsheets Skills}\label{spreadsheets-skills}

\part{Statistics}\label{part-statistics}

\chapter{Probability}\label{probability}

\chapter{Statistical Literacy}\label{statistical-literacy}

\chapter{Statistical Diagrams}\label{statistical-diagrams}

\chapter{Descriptive Statistics}\label{descriptive-statistics}

\chapter{Hypothesis Testing}\label{hypothesis-testing}

\chapter{Confidence Intervals}\label{confidence-intervals}

\part{Project}\label{part-project}

\chapter*{Overview}\label{overview}
\addcontentsline{toc}{chapter}{Overview}

\part{Finance}\label{part-finance}

\chapter{Financial Products}\label{financial-products}

\chapter{Interest}\label{interest}

\chapter{Salaries and Pensions}\label{salaries-and-pensions}

\part{Planning}\label{part-planning}

\chapter{Project Planning}\label{project-planning}

\chapter{Expectated Value}\label{expectated-value}

\part{Revision}\label{part-revision}

\chapter*{September Assessment}\label{september-assessment}
\addcontentsline{toc}{chapter}{September Assessment}

\bibliography{book.bib,packages.bib}

\end{document}
